\section{Explaining the recovery curve and specifying its scope}
\label{sec:blind}

The earlier internal analysis compared estimated intrinsic effects with the
requested reference in Equation~\ref{eq:requested}. Ratios were 1.00--1.07
where the rule returned an estimate on reference-only draws and 0.86--1.18 on
pooled draws, with shrinking dispersion at larger counts
(Appendix~\ref{app:generator}). The oracle estimator applies
Equation~\ref{eq:kitagawa} to the draw's own realised summaries. Its residual
against that empirical decomposition is zero across 1{,}711{,}200 diagnostic
replicates on two cohorts. The raw calculation also includes zero-cell cases;
those arithmetic outputs do not define an intrinsic effect in an absent
population (Appendix~\ref{app:generator}). On defined contrasts, the recovery
ratio therefore describes realised over requested reference. It does not
measure whether the reporting gate meets its interval requirements.

\paragraph{The general form.}
Let $D$ be a simulated draw requested at parameter $\theta$. Pre-specify an
observed-data reference $T(D)=h(S(D))$, where $S$ denotes summaries of the draw.
We call $T(D)$ the \emph{realised reference}; it need not be a realised causal
effect. If the estimate uses the same functional,
$\hat\theta(D)=h(S(D))$, then, for $\theta\ne0$,
\[
\frac{\hat\theta(D)}{\theta}=\frac{T(D)}{\theta},
\qquad r(D)=\hat\theta(D)-T(D)\equiv0.
\]
\begin{center}
\texttt{r = estimate[valid] - reference[valid]}\\
\texttt{max\_residual = np.max(np.abs(r))}
\end{center}
Here \texttt{valid} marks paired finite outputs; report excluded draws and do
not compute the maximum if none remain. This is a numerical discrepancy check, not a validity screen. Report the
raw maximum and outcome units; if using a closeness label, state both absolute
and relative tolerances and any sensitivity to that choice. No universal
tolerance separates correct from incorrect estimators. Only an algebraic
argument establishes an identity. Sharing summaries alone
is insufficient: different functions of the same summaries need not agree.

\paragraph{What the audit can establish.}
The recovery curve still measures sampling performance against $\theta$.
If estimate and reference are equivalent, their curves are redundant for
comparing those implementations. That is compatible with an accurate estimator
and a useful recovery experiment. The unsupported inference is from this
point-estimate comparison to interval coverage or an abstention rule, neither
of which the comparison evaluates. A zero residual therefore prompts inspection
of the claimed evidence; it is not a rejection criterion for an estimator.

\paragraph{A control separating equality from performance.}
To test this distinction directly, a separate simulation fixes the reference
before evaluating any estimator: for 200 independent observations
$Y_i\sim\mathcal N(\theta,1)$, set $T(D)=\bar Y$. We evaluate four procedures
at $\theta=0$ and $0.5$, with 2{,}000 replicates per effect. The baseline
reports the full-sample mean with a normal 95\% interval using the sample
standard deviation. A second procedure keeps that point estimate but quarters
the interval half-width. A third adds a known bias of $0.5$ to the baseline
point and interval. The fourth uses every other observation, selected before
seeing outcomes, with its own mean and standard error. This half-sample
estimate shares observations with the full-sample reference; it is not
independent of $T(D)$.

\begin{table}[htbp]
\centering\small
\caption{Equality and statistical performance in the clean control. All
2{,}000 outputs per row are finite. Residuals compare points with the
pre-specified full-sample mean; bias, RMSE and coverage compare with the
requested $\theta$. The nominal coverage is 95\%.}
\label{tab:cleancontrol}
\begin{tabular}{@{}clrrrr@{}}
\toprule
$\theta$ & procedure & $\max|r|$ & bias & RMSE & coverage \\
\midrule
$0$ & Full-sample mean & $0$ & $-0.0008$ & $0.0710$ & $94.60\%$ \\
 & Same point, narrow interval & $0$ & $-0.0008$ & $0.0710$ & $35.95\%$ \\
 & Mean plus $0.5$ & $0.5000$ & $0.4992$ & $0.5042$ & $0.00\%$ \\
 & Half-sample mean & $0.2883$ & $0.0004$ & $0.0994$ & $95.00\%$ \\
\midrule
$0.5$ & Full-sample mean & $0$ & $0.0001$ & $0.0710$ & $95.40\%$ \\
 & Same point, narrow interval & $0$ & $0.0001$ & $0.0710$ & $36.30\%$ \\
 & Mean plus $0.5$ & $0.5000$ & $0.5001$ & $0.5052$ & $0.00\%$ \\
 & Half-sample mean & $0.2686$ & $-0.0017$ & $0.1015$ & $94.20\%$ \\
\bottomrule
\end{tabular}
\end{table}

Table~\ref{tab:cleancontrol} puts adequate and inadequate intervals in the same
zero-residual class. Conversely, departure contains both the unbiased
half-sample mean and the deliberately biased estimate. Coverage Monte Carlo
standard errors are at most 1.08 percentage points. Neither residual class
certifies performance, even with a reference declared in advance. At the
null, differences remain well-defined although recovery ratios do not.


\paragraph{Report performance and equivalence separately.}
State the scientific target and whether it is fixed or draw-dependent, then
report target error and interval performance. If a realised reference is also
used, declare its functional and observation process before inspecting
residuals. Report attempted and valid pairs, maximum absolute residual in outcome units,
and both tolerances used for a numerical classification. Classify the result
as a proved identity on a stated domain, numerical agreement on the tested
draws, or an observed departure; only the first supports an identity claim. Establish identities algebraically
on a stated class of datasets; agreement on a restricted, noiseless design need
not extend beyond it. Evaluate abstention frequency and conditional interval
performance separately, using independent evaluation draws for any selected
rule. At null effects, report differences instead of undefined recovery ratios.
